\documentclass[a4paper,11pt]{article}
\usepackage[utf8]{inputenc}
\usepackage[spanish]{babel}
\usepackage{amsmath, amssymb, amsfonts}
\usepackage{geometry}
\geometry{a4paper, margin=2.5cm}
\usepackage{hyperref}
\usepackage{graphicx}
\usepackage{bookmark}

\title{\textbf{Matemática del Entorno AG-RACE y Dinámica del Agente}}
\author{Proyecto Jackpot Jockeys}
\date{\today}

\begin{document}

\maketitle

\section{Introducción}
Este documento describe el modelo matemático que rige las físicas del vehículo (\textit{Hovercraft}) y las observaciones que el agente de Inteligencia Artificial (IA) realiza dentro del entorno \textbf{AG-RACE}. El sistema está diseñado tanto para ser evaluado por heurísticas matemáticas (\textit{Pure Pursuit}) como para entrenar políticas de Aprendizaje por Refuerzo (RL).

\section{Dinámica del Vehículo (Modelo Cinemático)}
El vehículo se modela en un espacio bidimensional (2D) sujeto a aceleraciones lineales, fricción simulada ("drag") y atenuación de derrapes.

\subsection{Estado del Vehículo}
El estado dinámico del vehículo se define por su posición global $\mathbf{p} \in \mathbb{R}^2$, su velocidad lineal $\mathbf{v} \in \mathbb{R}^2$ y su rotación (orientación) $\theta \in [-\pi, \pi]$.

\subsection{Vectores Direccionales}
A partir del ángulo $\theta$, se derivan los vectores dirección frontal (forward) y derecho (right):
$$ \hat{\mathbf{f}} = \begin{bmatrix} \cos\theta \\ \sin\theta \end{bmatrix}, \quad \hat{\mathbf{r}} = \begin{bmatrix} -\sin\theta \\ \cos\theta \end{bmatrix} $$

\subsection{Entradas de Control (Acciones)}
El agente proporciona tres valores escalares continuos procesados por fotograma interactivo:
\begin{itemize}
    \item \textbf{Aceleración ($T$):} $T \in [0, 1]$
    \item \textbf{Freno ($B$):} $B \in [0, 1]$
    \item \textbf{Giro ($S$):} $S \in [-1, 1]$
\end{itemize}

\subsection{Actualización de Movimiento y Giro}
La orientación se actualiza de manera dependiente de la entrada de giro y la velocidad de giro constante $\omega_{turn}$:
$$ \theta_{t+\Delta t} = \theta_t + (S \cdot \omega_{turn} \cdot \Delta t) $$

La fuerza de empuje modifica la velocidad a través de la aceleración constante $a_{max}$ y un modificador de fricción ambiental $\mu_{fric}$:
$$ \mathbf{v}_{t+\Delta t} = \mathbf{v}_t + \left( \hat{\mathbf{f}} \cdot T \cdot a_{max} \cdot \mu_{fric} \cdot \Delta t \right) $$

Si se aplica freno ($B > 0$), la velocidad se reduce linealmente en contra de su propio vector de movimiento:
$$ \mathbf{v}_{t+\Delta t} = \mathbf{v}_{t} - \left( \frac{\mathbf{v}_{t}}{|\mathbf{v}_{t}|} \cdot \min\left( B \cdot a_{brake} \cdot \mu_{fric} \cdot \Delta t, |\mathbf{v}_{t}| \right) \right) $$

\subsection{Fricción Diferencial (Drifting)}
Para simular el efecto \textit{Hovercraft}, se descompone la velocidad en sus magnitudes frontal ($v_f$) y lateral ($v_l$):
$$ v_f = \mathbf{v} \cdot \hat{\mathbf{f}} \quad \text{y} \quad v_l = \mathbf{v} \cdot \hat{\mathbf{r}} $$
La velocidad lateral de derrape se atenúa mediante un coeficiente de agarre lateral $K_{grip}$ sujeto a deterioro ambiental $\mu_{drift}$:
$$ v_{l, new} = v_l - \text{sgn}(v_l) \cdot \min(|v_l|, K_{grip} \cdot \mu_{drift} \cdot |v_l| \cdot \Delta t) $$
De igual modo, la velocidad frontal sufre atenuación por resistencia al aire (drag) $C_d$:
$$ v_{f, new} = v_f - \text{sgn}(v_f) \cdot \min(|v_f|, C_d \cdot \mu_{fric} \cdot |v_f| \cdot \Delta t) $$

La velocidad final compuesta y re-ensamblada sujeta a los límites físicos $V_{max}$ es:
$$ \mathbf{v}_{final} = v_{f, new}\hat{\mathbf{f}} + v_{l, new}\hat{\mathbf{r}} $$
$$ \mathbf{v}_{t+\Delta t} = \min( |\mathbf{v}_{final}|, V_{max} ) \frac{\mathbf{v}_{final}}{|\mathbf{v}_{final}|} $$

\section{Percepción del Agente (Políticas de RL)}
El entorno es formulado como un Proceso de Decisión de Markov (MDP) parcialmente observable.

\subsection{Proyección en la Pista (Espacio de Frenet)}
La curva de la pista central se evalúa como una función paramétrica $C(s) \in \mathbb{R}^2$ donde $s \in [0, L]$ asume $L$ como perímetro total de la pista.

Para cualquier posición global $\mathbf{p}$, existe el escalar $s_{closest}$ que minimiza la distancia eucíldea:
$$ d_{error}(\mathbf{p}) = \min_{s} \| \mathbf{p} - C(s) \| $$

Esta métrica $d_{error}$ es devuelta como magnitud absoluta penalizante ("off-track error"). Adviértase que, si el vehículo sobrepasa el ancho asimilado asintomáticamente de la pista ($W_{track}/2$), tal valor escala significativamente para penalizar severamente a la política RL.

\subsection{Error de Orientación (Heading Error)}
El error es obtenido a partir de una pre-sintonía de distancia proyectada $D_{lookahead}$ asumiendo:
$$ Target = C( (s_{closest} \cdot L + D_{lookahead}) \pmod L ) $$

El vector unitario óptimo de avance es $\hat{\mathbf{u}} = \frac{Target - \mathbf{p}}{|Target - \mathbf{p}|}$.
El error direccional alimentado a las observaciones se computa contra $\hat{\mathbf{f}}$:
$$ e_\theta = \arccos(\hat{\mathbf{u}} \cdot \hat{\mathbf{f}}) \quad \text{Normalizado bajo} \left[ \frac{e_\theta}{\pi} \right] $$

\subsection{Vector Estocástico}
Las observaciones recopiladas construyen un vector matricial $O_t \in \mathbb{R}^{20}$ combinando magnitudes intrínsecas del vehículo (velocidades relativas, variables de atascamiento), discrepancias espaciales (errores $e_\theta, d_{error}$) y perturbaciones del World Event Manager derivadas como ruido inyectado, determinando la dinámica de transición de las políticas de aprendizaje.

\end{document}
